\scanpage{73}
\mspage{79}

\uline{Proposition \unsure{2|3}} Soit $f:X\longrightarrow Y$ dans $\mathrm{Diag}$. Conditions équivalentes

a) Pour tout $Y''\xrightarrow{p}Y'\xrightarrow{\unsure{g|r}}Y$, d'où
\diagram{\begin{tikzcd}[ampersand replacement=\&]
X''\arrow[r,"q"]\arrow[d] \& X'\arrow[r]\arrow[d] \& X\arrow[d,"f"]\\
Y''\arrow[r,"p"] \arrow[ur,phantom,"\text{cart}" description] \& Y'\arrow[r]\arrow[ur,phantom,"\text{cart}" description] \& Y
\end{tikzcd}}
Si $p$ est une $\mathbb D$-équivalence, $q$ aussi.

\margin{$\longrightarrow$ il faut supposer $Y',Y''$ (Hot)-asphériques cf page \unsure{96|95}, et supposer vérifié l'axiome $\mathrm{Der}\,3\,\unsure{a'|a})$ de la page \unsure{97}.}

% slanted
\margin{\uuline{$b)\Rightarrow a)$ \unsure{canular}} cf annotation marginale cf contre-exemple page 89. Mais ça marche \unsure{p.} pour les dérivateurs $\mathbb D_{\unsure{i|\delta}}$ $(i\in\mathbb N)$ associés aux ``localisateurs fondamentaux'' $W_i$, i.e. la théorie des \unsure{$i$|$\delta$}-types d'homotopie}

% slanted
\margin{L'implication $b\Rightarrow a$ pas prouvée, seulement sous forme plus faible : a$'$) si $Y'',Y'$ \deleted{\unsure{sont}} $\mathbb D$-asphériques et 1-connexes, \deleted{\unsure{c'est alors que}} alors $X''\longrightarrow X'$ est une $\mathbb D$ équivalence. (Il suffit $=$ de \unsure{de le|le|la} supposer pour $Y'$, $Y''$ \deleted{\unsure{tel}} ayant des objets ``\unsure{finaux}'' (pour \unsure{avoir} qu'a$'$ $\Rightarrow$ b) ou ``ayant des objets initiaux'' pour impliquer \unsure{a$'$)$\Rightarrow$b$'$)}}

b) \deleted{\unsure{Donc}} $f$ est une $\mathbb D$-fibration cohomologique \unsure{univ\textsuperscript{lle}}, i.e. elle est une fibration \inserted{coh.}, et le \unsure{reste} par tout changement de base.

b$'$) Condition duale de b) : \deleted{\unsure{$\mathbb D$}} $f$ est une $\mathbb D$-fibration homologique universelle.

$a)\Longrightarrow b)$ Supposons a), Soit $Y'\to Y$ chgt de base, d'où
\diagram{\begin{tikzcd}[ampersand replacement=\&]
X'\arrow[r]\arrow[d,"f'"']\& X\arrow[d,"f"]\\
Y'\arrow[r]\&Y
\end{tikzcd}}
. Pour prouver que $f'$ est une $\mathbb D$-fib coh. on regarde au dessus de $Y'$ les foncteurs de la forme $\{y'\}\longrightarrow Y'_{/y'}$ $(y'\in\mathrm{Ob}\,Y')$ et $Y'_{/y'_1}\longrightarrow Y'_{/y'_0}$ (si $u':y'_1\unsure{-|\longrightarrow}y'_0$ flèche dans $Y'$), et on veut que les flèches \deleted{déduites} de $\mathrm{Cat}_{/Y'}$ déduites de celles-là par le chgt de base $X'\xrightarrow{f'}Y'$ soient des $\mathbb D$-équiv.
% archive sheet number: 72
